% Bagian: first-order-logic
% Bab: proof-systems
% Subbagian: seqeunt-calculus

\documentclass[../../../include/open-logic-section]{subfiles}

\begin{document}

\iftag{FOL}
      {\olfileid[id]{fol}{prf}{seq}}
      {\olfileid[id]{pl}{prf}{seq}}

\olsection{Kalkulus Sekuen}

Sementara banyak sistem !!{derivation} bekerja dengan susunan
!!{sentence}s, kalkulus sekuen bekerja dengan \emph{sekuen}. Sekuen
adalah ekspresi berbentuk
\[
!A_1, \dots, !A_m \Sequent !B_1, \dots, !B_n,
\]
yaitu sepasang barisan !!{sentence}s yang dipisahkan oleh simbol
sekuen~$\Sequent$. Masing-masing barisan itu dapat kosong. !!^a{derivation}
dalam kalkulus sekuen merupakan pohon sekuen; sekuen paling atas
memiliki bentuk khusus (disebut ``sekuen awal'' atau ``aksioma''), dan
setiap sekuen lainnya mengikuti dari sekuen-sekuen tepat di atasnya
melalui salah satu aturan inferensi. Aturan inferensi dapat
memanipulasi !!{sentence}s dalam sekuen (menambahkan, menghapus, atau
menyusun ulang kalimat-kalimat itu di sisi kiri ataupun kanan), atau
memperkenalkan !!{formula} kompleks dalam kesimpulan aturan tersebut.
Sebagai contoh, aturan $\LeftR{\land}$ memungkinkan inferensi dari $!A,
\Gamma \Sequent \Delta$ ke $!A \land !B, \Gamma \Sequent \Delta$, dan
aturan $\RightR{\lif}$ memungkinkan inferensi dari $!A, \Gamma \Sequent
\Delta, !B$ ke $\Gamma \Sequent \Delta, !A \lif !B$, untuk sebarang
$\Gamma$, $\Delta$, $!A$, dan~$!B$. (Secara khusus, $\Gamma$ dan~$\Delta$
dapat kosong.)

Relasi $\Proves$ berdasarkan kalkulus sekuen didefinisikan sebagai
berikut: $\Gamma \Proves !A$ jika dan hanya jika terdapat suatu barisan
$\Gamma_0$ sedemikian sehingga setiap $!A$ dalam $\Gamma_0$ berada
dalam~$\Gamma$ dan terdapat suatu !!{derivation} dengan
sekuen~$\Gamma_0 \Sequent !A$ pada akarnya. $!A$ merupakan teorema dalam
kalkulus sekuen jika sekuen~$\Sequent !A$ memiliki !!a{derivation}.
Sebagai contoh, berikut adalah !!a{derivation} yang menunjukkan bahwa
$\Proves (!A \land !B) \lif !A$:
\begin{prooftree}
  \Axiom$!A \fCenter !A$
  \RightLabel{\LeftR{\land}}
  \UnaryInf$!A \land !B \fCenter !A$
  \RightLabel{\RightR{\lif}}
  \UnaryInf$\fCenter (!A \land !B) \lif !A$
\end{prooftree}

Himpunan $\Gamma$ tidak konsisten dalam kalkulus sekuen jika terdapat
!!a{derivation} dari $\Gamma_0 \Sequent$ (dengan setiap $!A \in \Gamma_0$
berada dalam~$\Gamma$ dan sisi kanan sekuen kosong). Dengan menggunakan
aturan \RightR{\Weakening}, setiap !!{sentence} dapat diturunkan dari
himpunan yang tidak konsisten.

Kalkulus sekuen diciptakan pada tahun 1930-an oleh Gerhard Gentzen.
Karena rancangannya sistematis dan simetris, kalkulus ini merupakan
formalisme yang sangat berguna untuk mengembangkan teori tentang
!!{derivation}s. Menemukan !!{derivation}s dalam kalkulus sekuen relatif
mudah, tetapi !!{derivation}s tersebut sering kali sulit dibaca dan
kaitannya dengan pembuktian kadang-kadang tidak mudah dilihat. Namun,
kalkulus sekuen telah terbukti sebagai pendekatan yang sangat elegan
terhadap sistem !!{derivation}, dan banyak logika memiliki sistem
kalkulus sekuen.

\end{document}
